\section{Agent report: relativistic Coulomb, polarized vacua, and low-energy limits}
\label{app-rel:report}

This report records the hypotheses and limitations of the relativistic Coulomb,
Bogoliubov--Dirac--Fock (BDF), and nonrelativistic-limit results used in this
paper. Statements attributed to the literature are separated from elementary
deductions and conditional estimates. None of the results reviewed below is a
stability theorem for the complete electron--positron--photon Fock-space model
with arbitrary point-nuclear configurations.

\subsection{The Lieb--Yau theorem and the relevant collapse mechanism}
\label{app-rel:lieb-yau}

Lieb and Yau \cite{LY1988} consider fixed-number fermionic electrons with
spin degeneracy $q$, fixed point nuclei, Coulomb interactions, and kinetic
energy $\sqrt{p^2+m^2}-m$ in units $\hbar=c=1$. Their stability conditions are
\begin{equation}
 Z_{\max}\alpha\le \frac{2}{\pi},\qquad
 \alpha<\frac{1}{47q}.
 \label{app-rel:ly-conditions}
\end{equation}
For physical electron spin, $q=2$, these inequalities include
$\alpha\simeq1/137$ and $Z_{\max}=26$. The model has no positrons, quantized
photons, or pair creation. The theorem therefore supports the electrostatic
part of the research program, but supplies no estimate for the omitted
operators.

The one-center critical coupling is $Z\alpha=2/\pi$. The paper also
establishes collective instability at sufficiently large $\alpha$, even for
arbitrarily small positive nuclear charge. Bounds on individual $Z\alpha$
are consequently not the only issue in relativistic many-body stability.

\paragraph{A scaling check.}
The sharp Coulomb form inequality on $\mathbb R^3$ is
$|p|\ge 2/(\pi|x|)$; its application and sharp threshold are part of
\cite{LY1988}. If $\nu>2/\pi$, sharpness gives a normalized form-domain
vector $f$ for which
\[
 c_f:=\langle f,|p|f\rangle-
 \nu\langle f,|x|^{-1}f\rangle<0.
\]
For $f_s(x)=s^{3/2}f(sx)$, functional calculus and
$\sqrt{a^2+m^2}-m\le a$ for $a\ge0$ give
\begin{equation}
 \left\langle f_s,
 \left(\sqrt{p^2+m^2}-m-\frac{\nu}{|x|}\right)f_s
 \right\rangle\le s c_f\longrightarrow-\infty.
 \label{app-rel:collapse-scaling}
\end{equation}
This is an actual variational collapse argument for the uncut
pseudorelativistic operator. It does not apply at fixed fermion momentum
cutoff: dilation eventually takes $f_s$ outside the cutoff space. Nor is
$2/\pi$ the self-adjointness threshold of the Dirac--Coulomb operator.

\subsection{What polarized-vacuum theorems establish}
\label{app-rel:bdf}

The initial Hainzl--Lewin--S\'er\'e construction
\cite{HLS2004} obtains a stable polarized vacuum through a fixed-point
argument, with smallness restrictions on the field, coupling, and cutoff.
The later variational result \cite{HLS2005}, Theorem~1 and equation~(29),
allows every finite cutoff and external density $n$ of finite Coulomb energy.
For $0\le\alpha<4/\pi$, it gives a minimizer and
\begin{align}
 \mathcal E^{n}_{\mathrm{BDF}}(Q)
 &\ge
 \left(1-\frac{\alpha\pi}{4}\right)
       \operatorname{Tr}_{P_-^0}(D^0Q)
 +\frac{\alpha}{2}D(\rho_Q-n,\rho_Q-n)
 -\frac{\alpha}{2}D(n,n),
 \label{app-rel:bdf-bound}\\
 D(f,g)&:=\iint_{\mathbb R^3\times\mathbb R^3}
             \frac{f(x)g(y)}{|x-y|}\,dx\,dy.
 \nonumber
\end{align}
Here $Q$ is the change in the one-body density matrix relative to the free
sea, and the generalized kinetic trace is nonnegative. BDF is a
mean-field, no-photon theory; its quasifree states describe polarization and
pair degrees of freedom, but not all correlated Fock-space states. Literal
delta nuclei fail the hypothesis $D(n,n)<\infty$.

\paragraph{Two logical checks on transferring this theorem.}
First, if $\mathcal S$ is a restricted variational class in a Hilbert space
$\mathcal H$, then
\begin{equation}
 \inf_{\substack{\Psi\in\mathcal H\\\|\Psi\|=1}}
 \langle\Psi,H\Psi\rangle
 \le
 \inf_{\Psi\in\mathcal S}\langle\Psi,H\Psi\rangle.
 \label{app-rel:variational-direction}
\end{equation}
A lower bound on the right side does not lower-bound the left side. Thus
even a BDF theorem incorporating arbitrarily many pairs cannot by itself
establish stability against arbitrary correlated pair states.

Second, a finite fermion momentum cutoff makes the external coupling to a
point source well-defined. In a box, every fermion bilinear has Fourier
support among differences of allowed momenta, hence within $|k|\le2\Lambda_f$.
If a distribution $n_{\Lambda_f}$ has the same Fourier coefficients as
$\sum_i Z_i\delta_{R_i}$ on this set, it produces exactly the same
external coupling. This algebraic equivalence permits a finite-Coulomb-energy
representative for that coupling. It does not identify the original
nucleus--nucleus energy with $D(n_{\Lambda_f},n_{\Lambda_f})/2$.

In particular, the last term of \eqref{app-rel:bdf-bound} is generally
configuration dependent and includes nuclear cross terms. Obtaining a
bound proportional to $K$ requires comparison with the prescribed nuclear
repulsion and control of the remaining self terms. Treating the filtered
representative as a new nuclear charge density without stating this change
would obscure the point-source assumption.

\subsection{Dirac--Coulomb self-adjointness is a different question}
\label{app-rel:dirac-coulomb}

For $D^0-\nu/|x|$, essential self-adjointness on
$C_c^\infty(\mathbb R^3\setminus\{0\};\mathbb C^4)$ holds for
$0\le\nu\le\sqrt3/2$, including the endpoint. The domain is $H^1$ for
strict $\nu<\sqrt3/2$. For $\sqrt3/2<\nu<1$, there is a unique
distinguished extension with domain in $H^{1/2}$. At $\nu=1$, a limiting
distinguished extension exists but its domain is not contained in
$H^{1/2}$ \cite[Section~2]{ELS2021}.

For a finite signed measure $\mu$ satisfying
$|\mu(\{R\})|<1$ at every point, \cite[Theorem~1]{ELS2021} constructs the
unique distinguished extension of $D^0-\mu*|x|^{-1}$ with domain in
$H^{1/2}$, obtained by norm-resolvent truncation of the singular potential.
Its essential spectrum is
\begin{equation}
 (-\infty,-1]\cup[1,\infty).
 \label{app-rel:dirac-essential-spectrum}
\end{equation}
This theorem does not assert semiboundedness. Distinct nuclei with
$Z\le26$ meet the individual atomic-weight condition at physical $\alpha$;
coincident sources instead have the sum of their charges as atomic weight.
Multicenter spectral gaps require separate analysis.

\paragraph{Why the negative sea must be treated correctly.}
Equation~\eqref{app-rel:dirac-essential-spectrum} prevents interpreting the
unprojected one-particle Dirac operator as a positive electron kinetic
energy. Already in the free problem, let $\widehat f_n$ be normalized,
supported in a fixed-radius ball centered at $ne_1$, and valued in the
negative spectral subspace of $D^0(p)$. Then
\[
 \langle f_n,D^0 f_n\rangle
 =-\int\sqrt{|p|^2+m^2}\,|\widehat f_n(p)|^2\,dp
 \longrightarrow-\infty.
\]
The electron--positron Fock construction instead fills the reference sea
and assigns positive energy to both particles and holes. This removes the
displayed free instability. It does not prove a bound for the interactions
or justify dropping pair terms.

\subsection{Established nonrelativistic limits and their scope}
\label{app-rel:nr-limits}

Hainzl--Lewin--S\'er\'e \cite{HLS2008} establish existence of charged BDF
minimizers under an HVZ-type binding condition. Their Theorem~3 considers
$\nu\in\mathcal C\cap L^1$, $\nu\ge0$, $Z=\int\nu$, fixed positive integer
$N$ satisfying $Z>N-1$, and cutoff $\Lambda=c\Lambda_0$. For sufficiently
large $c$, minimizers exist and
\begin{equation}
 \mathcal E_c^\nu(N)-N g_0(0)
       \longrightarrow E_{\mathrm{HF}}^\nu(N),
 \label{app-rel:bdf-nr-limit}
\end{equation}
with subsequential $H^1$ convergence of their electronic orbitals to a
Hartree--Fock minimizer. Here $g_0(0)$ is the threshold of the
self-consistent free mean-field Dirac operator. This is a fixed-$N,Z$
mean-field limit, not a macroscopic error estimate for full QED.

Arai \cite{Arai2003} proves a nonrelativistic Dirac--Maxwell limit leading
to a spin-$1/2$ Pauli--Fierz realization. The starting model contains one
Dirac particle and photon Fock space; it does not contain a quantized
Dirac field with variable electron and positron numbers. The theorem
therefore verifies a restricted relativistic-to-nonrelativistic connection,
without supplying the pair-elimination argument required here.

\paragraph{An elementary kinetic error bound.}
For $x\ge0$, convexity of $(1+x)^{-1/2}$ or direct differentiation gives
\[
 0\le 1+\frac{x}{2}-\sqrt{1+x}\le\frac{x^2}{8}.
\]
Consequently functional calculus yields
\begin{equation}
 0\le\frac{p^2}{2m}
 -\left(\sqrt{m^2c^4+c^2p^2}-mc^2\right)
 \le\frac{p^4}{8m^3c^2}.
 \label{app-rel:kinetic-error}
\end{equation}
On a sector of $N$ particles whose momenta are bounded by $p_*$, the
sum of these errors is at most $Np_*^4/(8m^3c^2)$. This proves an extensive
free-kinetic approximation under an explicit low-momentum assumption.
It neither controls the electromagnetic interaction nor establishes that
minimizers of the full theory lie in that sector.

\subsection{A conditional Feshbach estimate, with proof}
\label{app-rel:feshbach}

Let $P$ project onto a chosen no-positron subspace within a fixed charge
sector, and put $Q=1-P$. Suppose the compressions and off-diagonal maps
admit compatible self-adjoint or closed-form realizations. Write
$A=PHP$, $D=QHQ$, and $B=QHP$. Fix real $z$ and assume
\begin{equation}
 D-z\ge gQ,\qquad g>0.
 \label{app-rel:gap-assumption}
\end{equation}
On vectors in the domain of $B$, spectral calculus gives the rigorous form
inequality
\begin{equation}
 0\le B^*(D-z)^{-1}B\le g^{-1}B^*B.
 \label{app-rel:feshbach-bound}
\end{equation}
Indeed $0\le(D-z)^{-1}\le g^{-1}Q$, so application to $B\psi$ proves
\eqref{app-rel:feshbach-bound}. The effective operator at spectral
parameter $z$ is
\begin{equation}
 H_{\mathrm{eff}}(z)=A-B^*(D-z)^{-1}B.
 \label{app-rel:feshbach-operator}
\end{equation}
For vectors on a common core, completion of the quadratic form gives
\begin{align}
 \langle(\psi,\chi),(H-z)(\psi,\chi)\rangle
 &=\left\|(D-z)^{1/2}\chi+(D-z)^{-1/2}B\psi\right\|^2
 \nonumber\\
 &\quad+\langle\psi,
       (A-z-B^*(D-z)^{-1}B)\psi\rangle.
 \label{app-rel:feshbach-square}
\end{align}
The identity extends whenever the displayed forms are closed. Thus, under
these hypotheses, positivity of $H_{\mathrm{eff}}(z)-z$ implies
positivity of $H-z$. No commutation of $B$ with $D$ is used.

The difficult hypotheses have not disappeared. One needs
\eqref{app-rel:gap-assumption} uniformly over the configurations under
consideration, and a form estimate for the correction whose error is
proportional to system size. Photon-dependent $B$ need not be a bounded
operator on the entire photon Fock space. An operator norm estimate after
an additional finite-energy projection is legitimate, but its dependence
on that projection must be tracked. A bound growing as $K^2$, $V^2$, or an
uncontrolled cutoff-dependent mode count cannot transfer a uniform
$-CK$ stability theorem.

The phrase ``below pair threshold'' is therefore insufficient on its own.
The relevant threshold is that of the interacting complementary
compression $QHQ$, with the external field and reference vacuum fixed.
Soft photons also mean that the no-positron sector has no isolated
finite-dimensional band in infinite volume. They do not by themselves
eliminate the free pair threshold; they invalidate an unqualified
finite-dimensional gapped-band argument.

\subsection{Consequences for the proposed stability proof}
\label{app-rel:consequences}

The verified results establish substantial stability, binding, and
nonrelativistic-limit statements in their stated models. They leave three
distinct steps for full QED: a lower bound on all correlated pair states,
a nonperturbative estimate for the transverse photon interaction, and a
vacuum-relative comparison uniform in the number and placement of nuclei.
The conditional formula above identifies one possible route to an
effective description, but its gap and extensive-error assumptions require
new proof. Stability of an effective Hamiltonian, without these estimates
and control of the complementary subspace, does not prove stability of
the underlying relativistic theory.
